%%FILE blueprint/src/macros/macros_common.tex
\newcommand{\co}{\!:}
\newcommand{\define}[1]{\index{#1}\emph{#1}} % définition du mot #1
\newcommand{\subdefine}[2]{\index{#1!#2}\emph{#2}} % définition du mot #2 comme sous-entrée de #1
\newcommand{\definex}[2]{\index{#2}\emph{#1}} % affiche #1 index #2
\newcommand{\subdefinex}[3]{\index{#2!#3}\emph{#1}} % affiche #1 index #3 comme sous-entrée de #2

\newcommand{\image}[1]{{\centering\includegraphics{images/#1}\par}}

\newcommand{\A}{\mathcal{A}} % un atlas
\newcommand{\E}{\mathcal{E}} % un espace affine
\newcommand{\Cinf}{C^\infty}
\newcommand{\Rel}{\mathcal{R}} % une relation différentielle

\newcommand{\rst}[2]{#1|_{#2}}
\newcommand{\bigO}[1]{O\left(#1\right)}

\DeclareMathOperator{\Diff}{Diff}
\DeclareMathOperator{\Loop}{\mathcal{L}}
\DeclareMathOperator{\Bij}{Bij}
\DeclareMathOperator{\Aut}{Aut}
\DeclareMathOperator{\Id}{Id}
\DeclareMathOperator{\rk}{rk}
\DeclareMathOperator{\GL}{GL}
\DeclareMathOperator{\SO}{SO}
\DeclareMathOperator{\SU}{SU}
\DeclareMathOperator{\Lin}{L}
\DeclareMathOperator{\Hom}{Hom}
\DeclareMathOperator{\Bsym}{B_\mathrm{s}}
\DeclareMathOperator{\Hess}{Hess}
\DeclareMathOperator{\im}{im}
\DeclareMathOperator{\coker}{coker}
\DeclareMathOperator{\supp}{supp} % support
\DeclareMathOperator{\ind}{ind} % indice d'un point critique
\DeclareMathOperator{\sgn}{ε} % signature d'une permutation
\DeclareMathOperator{\Alt}{Alt} % antisymétrisation
\DeclareMathOperator{\Div}{div} % divergence
\DeclareMathOperator{\tr}{tr} % trace
\DeclareMathOperator{\pr}{pr} % une projection
\DeclareMathOperator{\Int}{Int} % intérieur
\DeclareMathOperator{\Conv}{Conv} % enveloppe convexe
\DeclareMathOperator{\IntConv}{IntConv} % intérieur enveloppe convexe
\DeclareMathOperator{\Conn}{Conn} % connected component
\DeclareMathOperator{\Span}{Span} % sev engendré
\DeclareMathOperator{\Gr}{Gr} % Grassmanienne
\DeclareMathOperator{\CP}{CP} % Corrugation process
\DeclareMathOperator{\Sec}{Sec} % Sections
\DeclareMathOperator{\Hol}{Hol} % Holonomic sections
\DeclareMathOperator{\codim}{codim} % codimension

% Cohomologie de de Rham, version avec indice dR ou non
%\newcommand{\HdR}[1][*]{H^{#1}_{\!\scriptscriptstyle d\!R}}
%\newcommand{\HdRc}[1][*]{H^{#1}_{\!\scriptscriptstyle d\!R, c}}
\newcommand{\HdR}[1][*]{H^{#1}}
\newcommand{\HdRc}[1][*]{H^{#1}_c}

\newcommand{\transp}[1]{#1^T}

\newcommand{\eE}{\mathcal{E}}
\newcommand{\fF}{\mathcal{F}}
\newcommand{\tT}{\mathcal{T}}

\newcommand{\X}{\mathfrak{X}} % champs de vecteurs

\newcommand{\F}{\mathcal{F}} % non-holonomic section

\newcommand{\pt}{\mathrm{pt}} % point

\DeclareMathOperator{\Op}{Op} % non-specified open neighborhood
\DeclareMathOperator{\dCzero}{d_{C^0}} % C^0 distance
\DeclareMathOperator{\dCone}{d_{C^1}} % C^1 distance
\DeclareMathOperator{\bs}{bs} % base map of a jet bundle section


\newcommand{\bbn}{\mathbb{N}}
\newcommand{\bbz}{\mathbb{Z}}
%\newcommand*{\bbe}{\mathbb{E}}
\newcommand{\abs}[1]{\left\lvert #1\right\rvert}
\newcommand{\Abs}[1]{\lvert #1\rvert}
\newcommand{\brac}[1]{\left( #1\right)}

%%FILE blueprint/src/macros/macros_print.tex
% Déclaration de théorème avec thmtools
\declaretheorem[numberwithin=chapter]{theorem}
\declaretheorem[sibling=theorem]{proposition}
\declaretheorem[sibling=theorem]{corollary}
\declaretheorem[sibling=theorem]{remark}
\declaretheorem[sibling=theorem]{lemma}
\declaretheorem[sibling=theorem]{definition}
\declaretheorem[sibling=theorem]{example}

\declaretheorem[numbered=no, name=Theorem]{theorem-intro}
\declaretheorem[numbered=no, name=Proposition]{proposition-intro}
\declaretheorem[numbered=no, name=Corollary]{corollary-intro}
\declaretheorem[numbered=no, name=Lemma]{lemma-intro}
\declaretheorem[numbered=no, name=Definition]{definition-intro}

% Cache les commandes plastex
\newcommand{\uses}[1]{}
\newcommand{\proves}[1]{}
\newcommand{\lean}[1]{}
\newcommand{\leanok}{}

%\newcommand{\D}[1]{\gls{D}[#1]}
\newcommand{\D}[1]{\mathrm{Diff}(#1)}

\newenvironment{diagram}[1][]{\[\begin{tikzcd}[#1]}{\end{tikzcd}\]}
\newenvironment{dessin}[1][]{\[\begin{tikzpicture}[#1]}{\end{tikzpicture}\]}

% Bugfix :
% http://tex.stackexchange.com/questions/262617/issue-with-xrightarrow
\makeatletter
  \patchcmd{\arrowfill@}{-7mu}{-14mu}{}{}
  \patchcmd{\arrowfill@}{-7mu}{-14mu}{}{}
  \patchcmd{\arrowfill@}{-2mu}{-4mu}{}{}
  \patchcmd{\arrowfill@}{-2mu}{-4mu}{}{}
\makeatother

\newcommand{\oiint}{\iint}

\definecolor{fond}{RGB}{255,255,255}

\newcommand{\N}[1]{\mathcal{N}_{\!\!#1}} % neighborhood

%%FILE blueprint/src/macros/macros_web.tex
\newtheorem{theorem}{Theorem}[chapter]
\newtheorem{definition}[theorem]{Definition}
\newtheorem{proposition}[theorem]{Proposition}
\newtheorem{lemma}[theorem]{Lemma}
\newtheorem{sublemma}[theorem]{Sub-lemma}
\newtheorem{corollary}[theorem]{Corollary}
\newtheorem{remark}[theorem]{Remark}
\newtheorem{example}[theorem]{Example}

\newtheorem*{theorem-intro}{Theorem}[chapter]
\newtheorem*{definition-intro}[theorem]{Definition}
\newtheorem*{proposition-intro}[theorem]{Proposition}
\newtheorem*{lemma-intro}[theorem]{Lemma}
\newtheorem*{corollary-intro}[theorem]{Corollary}

\newcommand{\D}[1]{\mathrm{Diff}(#1)}

\newenvironment{diagram}[1][]{\begin{tikzcd}[#1]}{\end{tikzcd}}
\newenvironment{dessin}[1][]{\begin{tikzpicture}[#1]}{\end{tikzpicture}}
\renewcommand\qedhere{}

\newcommand{\intprod}{\,\lrcorner\,}

\newcommand{\N}[1]{\mathcal{N}_{#1}} % neighborhood
